% Part: incompleteness
% Chapter: arithmetization-syntax
% Section: introduction

\documentclass[../../../include/open-logic-section]{subfiles}

\begin{document}

\olfileid{inc}{art}{int}
\olsection{পরিচয়}

গণনসাধ্যতা ও যুক্তিবিদ্যার মধ্যে সংযোগ স্থাপন করতে হলে যুক্তিবিদ্যার বস্তুগুলি
(প্রতীক, পদ, !!{formula}s, !!{derivation}s), তাদের উপর ক্রিয়া, এবং তাদের ধর্ম
ও পারস্পরিক সম্বন্ধ নিয়ে এমনভাবে কথা বলতে হবে, যা গণনামূলকভাবে বিচার করা
যায়। প্রতীক, প্রতীকের অনুক্রম এবং সেগুলি থেকে গঠিত অন্যান্য বস্তুর উপর
গণনসাধ্য অপেক্ষক ও সম্বন্ধ বিবেচনা করে এটি সরাসরি করা যায়। যৌক্তিক
সংকেতবিন্যাসের সব বস্তু সসীম এবং প্রতীকের !!a{enumerable} সেট থেকে গঠিত বলে
গণনার কিছু মডেলে এটি সম্ভব। কিন্তু পুনরাবৃত্ত অপেক্ষকের মতো অন্য মডেলগুলি
কেবল সংখ্যা এবং সংখ্যার উপর সম্বন্ধ ও অপেক্ষকে সীমাবদ্ধ। তদুপরি, শেষ পর্যন্ত
আমরা কয়েকটি তত্ত্বের ভিতরেই সংকেতবিন্যাস নিয়ে কাজ করতে চাই—বিশেষ করে
পাটিগণিতের ভাষায় সূত্রায়িত তত্ত্বে। সে ক্ষেত্রে সংকেতবিন্যাসকে
\emph{পাটিগণিতায়িত} করা আবশ্যক; অর্থাৎ সংকেতবিন্যাসগত বস্তু, তাদের উপর ক্রিয়া
এবং তাদের সম্বন্ধকে যথাক্রমে সংখ্যা, পাটিগাণিতিক অপেক্ষক এবং পাটিগাণিতিক
সম্বন্ধ হিসেবে উপস্থাপন করতে হবে। লাইবনিজ পর্যন্ত প্রসারিত এই ধারণায়
সংকেতবিন্যাসগত বস্তুকে সংখ্যা দেওয়া হয়।

প্রতীকগুলিকে তাদের ``কোড'' হিসেবে সংখ্যা দেওয়া তুলনামূলকভাবে সহজ। কয়েকটি
প্রতীক কিছুটা অসুবিধা সৃষ্টি করে; যেমন, অসীমসংখ্যক !!{variable}s আছে, এমনকি
প্রত্যেক স্থানসংখ্যা~$n$-এরও অসীমসংখ্যক !!{function}s আছে। তবে অবশ্যই প্রতীকগুলিকে
এমন সুসংবদ্ধভাবে সংখ্যা দেওয়া যায় যাতে, যেমন, $\Obj v_2$ ও $\Obj v_3$-কে
আলাদা কোড দেওয়া হয়। প্রতীকের অনুক্রম (যেমন পদ ও !!{formula}s) সামলানো আরও
কঠিন। কিন্তু সংখ্যার অনুক্রমকে যদি সম্পূর্ণ পাটিগাণিতিক উপায়ে সামলাতে পারি
(যেমন অনুক্রমের মৌলিক সংখ্যার ঘাতভিত্তিক সংকেতায়ন দিয়ে), তবে আলাদা প্রতীকের
কোড থেকে প্রতীক-অনুক্রমের কোডে, এবং সেখান থেকে !!{formula}s-এর অনুক্রম বা
অন্য বিন্যাস—যেমন !!{derivation}s—এর কোডে যেতে পারি। এই বিস্তৃত সংকেতায়নকে
``গ্যোডেল সংখ্যায়ন'' বলা হয়। প্রত্যেক পদ, !!{formula} ও !!{derivation}-কে
একটি গ্যোডেল সংখ্যা দেওয়া হয়।

প্রতীক-অনুক্রমকে তার প্রতীকগুলির কোডের অনুক্রম হিসেবে সংকেতায়িত করে, এবং
অনুক্রমের এমন সংকেতায়ন বেছে নিয়ে যা গণনসাধ্য অপেক্ষক দিয়ে সামলানো যায়,
আমরা গ্যোডেল সংখ্যাকেও গণনসাধ্য অপেক্ষক দিয়ে সামলাতে পারি। বাস্তবে প্রাসঙ্গিক
সব অপেক্ষকই আদিম পুনরাবৃত্ত হবে। যেমন, কোনো অনুক্রমের দৈর্ঘ্য নির্ণয় এবং তার
কোড থেকে অনুক্রমের $i$-তম উপাদান নির্ণয়—দুটিই আদিম পুনরাবৃত্ত। অনুক্রমটির
কোড যদি, যেমন, কোনো !!a{formula}~$!A$-এর গ্যোডেল সংখ্যা হয়, তবে সঙ্গে সঙ্গে
দেখা যায় যে $!A$-এর গ্যোডেল সংখ্যা থেকে !!a{formula}-টির দৈর্ঘ্য এবং
!!a{formula}-টির $i$-তম প্রতীকের (কোড) নির্ণয় করা যায়। কোনো সংখ্যা সুগঠিত
পদের গ্যোডেল সংখ্যা কি না, অথবা সঠিক !!{derivation}-এর গ্যোডেল সংখ্যা কি
না—এই ধর্ম আদিম
পুনরাবৃত্ত প্রমাণ করা একটু কঠিন। তবুও তা সম্ভব, কারণ আমাদের প্রয়োজনীয়
অনুক্রমগুলি—পদ, !!{formula}s ও !!{derivation}s—আরোহমূলকভাবে সংজ্ঞায়িত।

উদাহরণ হিসেবে প্রতিস্থাপন ক্রিয়াটি বিবেচনা করো। $!A$ একটি সূত্র, $x$ একটি
চলরাশি এবং $t$ একটি পদ হলে $\Subst{!A}{t}{x}$ হলো $!A$-তে $x$-এর প্রতিটি
মুক্ত ঘটনাকে $t$ দিয়ে প্রতিস্থাপনের ফল। এখন ধরি, $!A$, $x$, $t$-কে যথাক্রমে
$k$, $l$ ও $m$ গ্যোডেল সংখ্যা দেওয়া হয়েছে। একই বিন্যাসে
$\Subst{!A}{t}{x}$-এরও একটি গ্যোডেল সংখ্যা, ধরা যাক $n$, পাওয়া যায়। $k$,
$l$ ও $m$-কে $n$-এ পাঠানো এই চিত্রণই প্রতিস্থাপন ক্রিয়ার পাটিগাণিতিক অনুরূপ।
প্রতিস্থাপন ক্রিয়া যখন $!A$, $x$, $t$-কে $\Subst{!A}{t}{x}$-এ পাঠায়, তখন
পাটিগণিতায়িত প্রতিস্থাপন অপেক্ষকটি গ্যোডেল সংখ্যা $k$, $l$, $m$-কে গ্যোডেল
সংখ্যা~$n$-এ পাঠায়। আমরা দেখব, এই অপেক্ষক আদিম পুনরাবৃত্ত।

সংকেতবিন্যাসের পাটিগণিতায়ন কেবল বিমূর্ত আগ্রহের বিষয় নয়। পাটিগণিতের ভাষার
মতো যে ভাষায় অন্য ভাষা ``নিয়ে কথা বলার'' নিজস্ব ব্যবস্থা নেই, সেই ভাষাতেও
জটিল অভিব্যক্তিগত ধর্ম সূত্রায়িত করা যায়—এটি প্রথমে এক গুরুত্বপূর্ণ ও
অনায়াসসাধ্য অন্তর্দৃষ্টি ছিল। এরপর স্বাভাবিক প্রশ্ন হলো, পেয়ানো পাটিগণিতের
মতো এই ভাষার কোনো তত্ত্ব তার নিজের ভাষা সম্বন্ধে কী \emph{প্রমাণ} করতে পারে
(যেমন, কোনো !!{sentence}s প্রমাণযোগ্য বা সত্য কি না)। এ থেকে গ্যোডেলের
অপ্রমাণযোগ্যতা-সংক্রান্ত এবং তারস্কির সত্যের অসংজ্ঞায়নীয়তা-সংক্রান্ত বিখ্যাত
সীমা-উপপাদ্যগুলিতে পৌঁছাই। সংকেতবিন্যাস পাটিগণিতায়িত করার কৌশলটি গণনসাধ্যতা
তত্ত্বের কয়েকটি গুরুত্বপূর্ণ ফল প্রমাণেও দরকার—যেমন তত্ত্বগুলির গণনক্ষমতা বা
গণনসাধ্যতার বিভিন্ন মডেলের পারস্পরিক সম্পর্ক। সংকেতবিন্যাসের পাটিগণিতায়ন
অন্যান্য বস্তু ও ধর্ম পাটিগণিতায়িত করারও আদর্শ দেয়। যেমন, কনফিগারেশন ও গণনা
(ধরা যাক, টুরিং যন্ত্রের) একইভাবে পাটিগণিতায়িত করা যায়। ফলে এক মডেলের
গণনা (যেমন টুরিং যন্ত্র) অন্য মডেলে (যেমন পুনরাবৃত্ত অপেক্ষক) অনুকরণ করা
সম্ভব হয়।

\end{document}
